\documentclass[12pt]{article}
\usepackage[utf8]{inputenc}
\usepackage[T1]{fontenc}
\usepackage{hyperref}
\title{San Felipe Otlaltepec Popoloca Language Resources for K-12 Curriculum Development}
\author{Compiled by Claude Sonnet 5.5 (claude-sonnet-5-5) for Andr\'as Kornai}
\date{}
\begin{document}
\maketitle

\section{Introduction}
San Felipe Otlaltepec Popoloca (ISO 639-3 \texttt{pow}; Otomanguean, Puebla, Mexico) is a threatened language; speaker estimates vary widely (about 7,000 in the reference description, about 3,000 in ASJP). This report rests on one web search and is not exhaustive.

\section{Primary Educational Resources}
No school materials were found.

\section{Linguistic and Reference Materials}
\begin{itemize}
\item Glottolog entry for San Felipe Otlaltepec Popoloca. \url{https://glottolog.org/resource/languoid/id/sanf1264}
\item ASJP wordlist page (cites SIL's Western Popoloca dictionary among sources). \url{https://asjp.clld.org/languages/SAN_FELIPE_OTLALTEPEC_POPOLOCA}
\item Wikitongues language page (EGIDS 7, shifting). \url{https://wikitongues.org/languages/pow/}
\end{itemize}

\section{Religious and Community Resources}
No scripture or community organisation was identified in the results; the Global Recordings page below lists recordings of religious content, if any.

\section{Digital and Audio Resources}
\begin{itemize}
\item Global Recordings Network, Popoloca San Felipe Otlaltepec. \url{https://globalrecordings.net/en/language/pow}
\item AILLA: Vocabulary, Otlaltepec Popoloca (M.\ Swanton with speaker Valeriano Aguilar Arrioja, 1999--2002). \url{https://islandora-ailla.lib.utexas.edu/islandora/object/ailla%3A134810}
\item AILLA: Narratives, Otlaltepec Popoloca (three collections). \url{https://islandora-ailla.lib.utexas.edu/islandora/object/ailla%3A134805}, \url{https://islandora-ailla.lib.utexas.edu/islandora/object/ailla%3A134806}, \url{https://islandora-ailla.lib.utexas.edu/islandora/object/ailla%3A134809}
\end{itemize}
Some AILLA items may require access permission.

\section{Curriculum Development Context}
The language is not known to be taught in schools, and documentation is mainly archival (AILLA recordings) with no found pedagogical materials. Any curriculum would require new work with community speakers, using the archive recordings as a starting point.
\end{document}